% Overleaf: Türkçe CV
% pdfLaTeX | ATS uyumlu
\documentclass[11pt,a4paper]{article}

\usepackage[utf8]{inputenc}
\usepackage[T1]{fontenc}
\usepackage[turkish]{babel}
\usepackage[margin=1.4cm]{geometry}
\usepackage{enumitem}
\usepackage{hyperref}
\usepackage{xcolor}

\hypersetup{colorlinks=true, urlcolor=black, linkcolor=black}
\pagestyle{empty}
\setlength{\parindent}{0pt}
\setlength{\parskip}{0pt}

\usepackage[scaled=0.95]{helvet}
\renewcommand{\familydefault}{\sfdefault}

\newcommand{\cvsection}[1]{%
  \vspace{1.2em}%
  {\large\textbf{\MakeUppercase{#1}}}\\[0.2em]
  \hrule
  \vspace{0.6em}%
}
\newcommand{\jobentry}[4]{%
  \textbf{#1} \hfill {\small #2}\\
  \textit{#3} \hfill {\small #4}\\[0.2em]
}

\begin{document}

\begin{center}
  {\LARGE \textbf{Ece Dalpolat}}\\[0.4em]
  {\large Yapay Zeka Mühendisi \& Veri Bilimci}\\[0.4em]
  {\small
  +90 552 360 9850 \,|\, ecedlplt9850@gmail.com \,|\, İstanbul\\[0.2em]
  \href{https://ecedalpolat.github.io/ece-dalpolat-portfolio/}{Portföy} \,|\,
  \href{https://github.com/EceDalpolat}{GitHub} \,|\,
  \href{https://www.linkedin.com/in/ece-dalpolat-0265a0221/}{LinkedIn}
  }
\end{center}

\vspace{-0.5em}
\cvsection{Özet}

Üretim ortamında çalışan yapay zeka sistemleri kuruyorum; demo değil, canlı B2B ürünler. Qkare'de kurumsal müşterilerin kullandığı sistemlerin ana mühendisiyim: dashboard besleyen dbt veri katmanı, psikometrik veri üzerinde soru yanıtlayan LangChain ve LangGraph ajanları, KPI ve yetkinlik öneren doğrulanmış JSON çıktılı FastAPI servisleri.

Günlük işim: SQL'de veri modelleme (109+ dbt modeli, native PostgreSQL satır düzeyi güvenlik, Superset), LLM hatları tasarlama (RAG, semantik önbellek, prompt injection koruması), ortak agent gateway arkasında mikroservis yayınlama, canlıya almadan önce LLM-as-judge ile kalite ölçme. Belirsiz ihtiyaçtan Docker/Kubernetes ve CI ile çok kiracılı, maliyet kontrollü üretime kadar uçtan uca sahipleniyorum.

\cvsection{Deneyim}

\jobentry{Yapay Zeka Mühendisi}{Haz 2025 -- Devam}{Qkare}{İstanbul}
\begin{itemize}[leftmargin=1.5em, nosep, topsep=2pt, itemsep=2pt]
  \item Ana mühendis: dbt katmanı, üretim AI ajanları/danışmanları, gateway ve güvenlik; 8+ canlı sistem.
\end{itemize}
\vspace{0.6em}

\jobentry{Öğretim Asistanı}{Oca 2025 -- Devam}{Miuul}{Uzaktan}
\begin{itemize}[leftmargin=1.5em, nosep, topsep=2pt, itemsep=2pt]
  \item Veri bilimi ve makine öğrenmesi dersleri; uygulamalı proje mentörlüğü ve değerlendirme.
\end{itemize}
\vspace{0.6em}

\jobentry{Lisansüstü Araştırma Görevlisi}{Oca 2024 -- Haz 2025}{Işık Üniversitesi}{İstanbul}
\begin{itemize}[leftmargin=1.5em, nosep, topsep=2pt, itemsep=2pt]
  \item Yazılım laboratuvarları; veri bilimi ve derin öğrenme projelerinde danışmanlık.
\end{itemize}
\vspace{0.6em}

\jobentry{Araştırma Veri Bilimci}{Oca 2024 -- Haz 2025}{Top4Honeychains}{}
\begin{itemize}[leftmargin=1.5em, nosep, topsep=2pt, itemsep=2pt]
  \item Bal değer zinciri ETL, TensorFlow/scikit-learn modelleri, REST API entegrasyonları.
\end{itemize}

\cvsection{Projeler}

\textbf{Analitik Platformu} \,|\, \textit{Qkare} \hfill Kas 2025 -- May 2026\\
Uçtan uca İK analitiği: PostgreSQL FDW, 109 dbt modeli, 56 mart, 13 AI ön hesaplama tablosu, native RLS, Superset. Sorumluluk: dbt katmanı.
\vspace{0.5em}

\textbf{Analyzer Platformu} \,|\, \textit{Qkare} \hfill Nis 2025 -- May 2026\\
Psikometrik mart'lar üzerinde LangChain ReAct ajanı; 4 katman güvenlik, JWT+RLS, Langfuse. Sorumluluk: güvenlik ve config bootstrap.
\vspace{0.5em}

\textbf{Coach Platformu} \,|\, \textit{Qkare} \hfill Eyl 2025 -- Kas 2025\\
Çok ajanlı koçluk (LangGraph, POML, OpenAI Realtime). Referee Agent (LLM-as-judge), ses katmanı, eval araçları.
\vspace{0.5em}

\textbf{Skills Advisor} \,|\, \textit{Qkare} \hfill Ağu 2025 -- May 2026\\
Kapalı sözlükten tam 12 yetkinlik boyutu; LangChain veya DSPy ensemble. Ana mühendis.
\vspace{0.5em}

\textbf{KPI Advisor} \,|\, \textit{Qkare} \hfill Nis 2026 -- May 2026\\
APQC PCF ile SMART KPI; lexical retrieval, Golden Dataset, LLM-as-judge hattı. Ana mühendis.
\vspace{0.5em}

\textbf{Agent API Gateway} \,|\, \textit{Qkare} \hfill Nis 2026 -- May 2026\\
FastAPI monorepo, birleşik invoke API, JWT, ajan başına Docker/CI. Platform sorumlusu.
\vspace{0.5em}

\textbf{QKare AI Asistan} \,|\, \textit{Qkare} \hfill Tem 2025\\
Kurumsal chatbot: kural, semantik önbellek (FAISS/Milvus), RAG, LLM; hibrit güvenlik. Tek yazar.
\vspace{0.5em}

\textbf{Davranışsal Kümeleme Ar-Ge} \,|\, \textit{Qkare} \hfill Ağu 2025\\
K-Means, DBSCAN, GMM; PCA/t-SNE; Analitik Platform mart'larına girdi. Tek yazar.
\vspace{0.5em}

\textbf{CVAnalyzerAI} \hfill 2025\\
MCP ile özgeçmiş analizi (FastAPI, DSPy, ANTLR). github.com/EceDalpolat/cvAnalizerAI
\vspace{0.5em}

\textbf{TOP4HONEYCHAINS} \hfill 2024--2025\\
Bal izlenebilirlik ML, REST API, zaman serisi tahmini.
\vspace{0.5em}

\textbf{Pazarlama Kampanyası Tahmini}\\
Banka kampanya modeli: SMOTE, XGBoost, Streamlit.
\vspace{0.5em}

\textbf{Polen Sınıflandırma}\\
CNN, \%90+ doğruluk; veri artırma.
\vspace{0.5em}

\textbf{Anten Davranış Analizi}\\
DeepLabCut ile keypoint takibi; hareket analizi.
\vspace{0.5em}

\textbf{Öğrenci Başarı Analizi}\\
Akademik özelliklerle Random Forest; Tkinter arayüzü.

\cvsection{Beceriler}

\textbf{AI/LLM:} LangChain, LangGraph, DSPy, RAG, OpenAI API, Pydantic, LLM-as-judge\\[0.2em]
\textbf{Veri:} dbt Core, PostgreSQL, Apache Superset, Pandas, SQL\\[0.2em]
\textbf{Backend:} FastAPI, httpx, WebSocket, JWT, MCP, REST\\[0.2em]
\textbf{Altyapı:} Docker, Kubernetes, GitLab CI/CD, Redis, OpenTelemetry, Langfuse\\[0.2em]
\textbf{Kod:} Python, SQL, JavaScript, Java, C\#

\cvsection{Eğitim}

\textbf{Yüksek Lisans, Bilgi Teknolojileri} \hfill 2022 -- Devam\\
Işık Üniversitesi --- GNO 3,57. Tez: arıcılık ve bal hasadında makine öğrenmesi otomasyonu.
\vspace{0.5em}

\textbf{Lisans, Yazılım Mühendisliği} \hfill 2018 -- 2023\\
Kırklareli Üniversitesi --- GNO 2,71

\cvsection{Sertifikalar}

Miuul Veri Bilimci Bootcamp; Miuul Makine Öğrenmesi; Miuul Üretken AI; Udemy LangChain/RAG; BTK Uygulamalı SQL; Huawei HCCDA-AI

\vspace{0.8em}
\begin{center}
  {\small Türkçe (ana dil) \,|\, İngilizce (B2)}
\end{center}

\end{document}
